\documentclass[10pt,a4paper]{amsart}
\usepackage[margin=19mm]{geometry}
\usepackage{amsmath,amssymb,mathtools}
\usepackage[scr=rsfs,cal=euler]{mathalfa}
\usepackage{xfrac}
\usepackage[unicode,hidelinks,bookmarks=false]{hyperref}
\newcommand{\ie}{i.e.\ }
\newcommand{\cf}{cf.\ }
\newcommand{\resp}{respectively\ }
\newcommand{\Subsection}{\subsection}
\providecommand{\nobreakdash}{\nobreak\hbox{-}}
\setlength{\emergencystretch}{2em}
\multlinegap0pt
\usepackage{bm}
\usepackage{bbm}
\usepackage{dsfont}
\usepackage{xspace}
\usepackage{cancel}
\usepackage{mathtools}
\usepackage{amsthm}
\makeatletter
\@ifundefined{thm}{\newtheorem{thm}{Theorem}}{}
\makeatother
\DeclareMathOperator{\Par}{\mathbb{Y}}
\DeclareMathOperator{\SYT}{\mathrm{SYT}}
\title[Non-archimedean Infinite Hecke Algebra]{Non-archimedean Infinite Hecke Algebra}
\author{Milo Bechtloff Weising}
\date{Source: arXiv:2512.01835v2 (CC BY 4.0)}
\begin{document}
\maketitle

\noindent\emph{Source and licence.} Non-archimedean Infinite Hecke Algebra. Version arXiv:2512.01835v2 (CC BY 4.0), distributed under the Creative Commons Attribution 4.0 International licence, as recorded in the arXiv licence metadata for this exact version and shown on the abstract page https://arxiv.org/abs/2512.01835v2. Full rights notice: licenses/arxiv-2512.01835-CC-BY-4.0.txt.\par\smallskip

\noindent\textit{Editorial scope.} Counted unit: the Thm whose source label is \texttt{rationality thm} in the article (source0.tex lines 923--925, source label \texttt{rationality thm}), delivered together with the article's own complete original proof of it (source0.tex lines 926--963, one \texttt{proof} environment of 38 lines). No mathematics is written, repaired or re-derived: statement and proof are the article's own text line for line. Required context retained uncounted: none. Every definition, display and label of those lines is retained unchanged. Omitted from this package and not redistributed: 1--922, 964--1000. Nothing inside a retained excerpt is omitted or altered: every retained body is delivered exactly as the article's lines are written, so no omission receipt of any kind accompanies this package. Citations: the retained excerpts carry no citation command at all, so no bibliography entry of the article is transcribed and no reference is redistributed. Reference adaptation: none was needed and none was made; every reference inside the delivered unit resolves inside it, so no unresolved reference or citation is produced. Printed slips kept literal: no printed word, symbol, hypothesis or number of the counted statement or of its proof was altered. Layout and class: the article's own class is \texttt{amsart} with packages including geometry, amssymb, amsmath, mathrsfs, amsfonts, hyperref, xcolor, bbm, tikz, array, multirow, ytableau, thmtools, thm-restate; the wrapper is the lane's pinned \texttt{amsart} editorial wrapper at 10pt on A4, which restates the theorem-like environments the retained text needs and adds the article's own macro definitions that the retained text needs (\texttt{\textbackslash Par}, \texttt{\textbackslash SYT}), with extra wrapper packages (bm, bbm, dsfont, xspace, cancel, mathtools). The delivered numbering is the unit's own. Assets: no retained span contains a figure environment, a graphics-inclusion command, a PDF-page inclusion, a table or a TikZ picture; no image file is copied into this package, the package declares no asset dependency and no image file is redistributed with it. Licence: version arXiv:2512.01835v2 of this article is distributed under the Creative Commons Attribution 4.0 International licence, as recorded in the arXiv licence metadata for this exact version and shown on the abstract page https://arxiv.org/abs/2512.01835v2. A full rights notice is delivered at licenses/arxiv-2512.01835-CC-BY-4.0.txt. Characters: every retained excerpt is pure ASCII, so the delivered engine, XeLaTeX with its default Latin Modern text font, reports no missing character.
\begin{thm}\label{rationality thm}
    For all $\lambda \in \Par$, $\Gamma_{\lambda}(1) \in k(t).$ 
\end{thm}

\begin{proof}
    Our strategy is to prove a recursive formula for $\Gamma_{\lambda}(1)$ which inductively implies that $\Gamma_{\lambda}(1) \in k(t)$. Let $\lambda \in \Par \setminus \{\emptyset\}.$ Recall that the diagram $\lambda^{(\infty)}$ contains a copy of the diagram $\lambda$ namely, $\lambda \equiv \lambda^{(\infty)}/\emptyset^{(\infty)}$. It is easy to see that for all $\tau \in \SYT_{\infty}(\lambda)$, the unique box in $\lambda^{(\infty)}$ labeled $\mathrm{rk}(\tau)$ by $\tau$ is a removable box in $\lambda$. For each removable box $\square_{0}$ of $\lambda$ define $S_{\square_{0}}$ as the set of $\tau \in \SYT_{\infty}(\lambda)$ such that $\tau(\square_0) = \mathrm{rk}(\tau).$ Necessarily,
    $$\SYT_{\infty}(\lambda) = \bigsqcup_{\square_0~ \text{removable}} S_{\square_0}$$ so 

    $$\Gamma_{\lambda}(1) = \sum_{\tau \in \SYT_{\infty}(\lambda)} \eta(\tau) = \sum_{\square_0~ \text{removable}} \sum_{\tau\in S_{\square_0}} \eta(\tau).$$ 

    Let $\square_0$ be a removable box of $\lambda$ considered in the diagram $\lambda^{(\infty)}$. Given a pair $(m,\tau')$ with $m\geq n_{\lambda}$ and $\tau' \in \SYT_{\infty}(\lambda \setminus \square_0)$ we construct a corresponding $\tau\in \SYT_{\infty}(\lambda)$ by defining $\tau(\square_0) := m$ and filling the boxes of $\lambda^{(\infty)}\setminus \{\square_0\} = (\lambda\setminus \{\square_0\})^{(\infty)}$ using the remaining labels $\{1,2,\dots\} \setminus \{m\}$ following the ordering determined by $\tau'.$ This is clearly a reversible construction so $S_{\square_0}$ is in bijection with $\{n_{\lambda},n_{\lambda}+1,\dots\} \times \SYT_{\infty}(\lambda \setminus \{\square_0\})$ via $(m,\tau') \mapsto \tau.$ It is straightforward casework to determine the following:

    $$\mathrm{Inv}(\tau) = \mathrm{Inv}(\tau') \sqcup \{ (\square_0,\square)| \square \in \lambda^{(\infty)}/\lambda^{(m)}\} \sqcup \{(\square_0,\square)|\square ~ \text{strictly to the right of} ~ \square_0 ~ \text{in} ~ \lambda^{(n_{\lambda})}\}.$$
    From this we directly compute,
    \begin{align*}
        \eta(\tau) &= \eta(\tau') t^{m-n_{\lambda}} \left(\frac{1-t^{\lambda_1-c(\square_0)-1}}{1-t^{\lambda_1-c(\square_0)+1}}\right)\cdots \left(\frac{1-t^{\lambda_1-c(\square_0)-1+m-n_{\lambda}-1}}{1-t^{\lambda_1-c(\square_0)+1+m-n_{\lambda}-1}}\right) \\
        & \times \prod_{\square_0 \rightarrow \square} t^{d(\square)+1}\left(\frac{1-t^{c(\square)-c(\square_0)-1}}{1-t^{c(\square)-c(\square_0)+1}}\right) \cdots \left(\frac{1-t^{c(\square)-c(\square_0)-1+d(\square)}}{1-t^{c(\square)-c(\square_0)+1+d(\square)}}\right)\\
    \end{align*}
    where $\square_0 \rightarrow \square$ ranges over boxes of $\lambda$ which are to the right of $\square_0$ and at the bottom of their column and $d(\square)$ is the number of boxes of $\lambda$ weakly above $\square$ in the same column (i.e., including $\square$ itself).

    Therefore, for all removable $\square_0,$
    \begin{align*}
        &\sum_{\tau\in S_{\square_0}} \eta(\tau) \\
        &=\sum_{m \geq n_{\lambda}} \sum_{\tau' \in \SYT_{\infty}(\lambda\setminus \{\square_0\})}\eta(\tau') t^{m-n_{\lambda}} \left(\frac{1-t^{\lambda_1-c(\square_0)-1}}{1-t^{\lambda_1-c(\square_0)+1}}\right)\cdots \left(\frac{1-t^{\lambda_1-c(\square_0)-1+m-n_{\lambda}}}{1-t^{\lambda_1-c(\square_0)+1+m-n_{\lambda}}}\right) \\
        & \times \prod_{\square_0 \rightarrow \square} t^{d(\square)+1}\left(\frac{1-t^{c(\square)-c(\square_0)-1}}{1-t^{c(\square)-c(\square_0)+1}}\right) \cdots \left(\frac{1-t^{c(\square)-c(\square_0)-1+d(\square)}}{1-t^{c(\square)-c(\square_0)+1+d(\square)}}\right)\\
        &= \Gamma_{\lambda \setminus \{\square_0\}}(1)\prod_{\square_0 \rightarrow \square} t^{d(\square)+1}\left(\frac{1-t^{c(\square)-c(\square_0)-1}}{1-t^{c(\square)-c(\square_0)+1}}\right) \cdots \left(\frac{1-t^{c(\square)-c(\square_0)-1+d(\square)}}{1-t^{c(\square)-c(\square_0)+1+d(\square)}}\right) \\
        & \times \sum_{m \geq n_{\lambda}} t^{m-n_{\lambda}} \left(\frac{1-t^{\lambda_1-c(\square_0)-1}}{1-t^{\lambda_1-c(\square_0)+1}}\right)\cdots \left(\frac{1-t^{\lambda_1-c(\square_0)-1+m-n_{\lambda}}}{1-t^{\lambda_1-c(\square_0)+1+m-n_{\lambda}}}\right).\\
    \end{align*}

Using a standard inductive argument on partial sums one may verify that
$$ \sum_{m \geq n_{\lambda}} t^{m-n_{\lambda}} \left(\frac{1-t^{\lambda_1-c(\square_0)-1}}{1-t^{\lambda_1-c(\square_0)+1}}\right)\cdots \left(\frac{1-t^{\lambda_1-c(\square_0)-1+m-n_{\lambda}}}{1-t^{\lambda_1-c(\square_0)+1+m-n_{\lambda}}}\right) = \frac{1-t^{\lambda_1-c(\square_0)}}{1-t}.$$ Note that here we are computing content values $c(\square)$ in $\lambda^{(n_{\lambda})}.$ If we instead compute the content of $\square \in \lambda$ in the diagram $\lambda$, we need to adjust accordingly. 

In the following formula, all relevant statistics are computed in the the diagram $\lambda$. Putting everything together we find:

\begin{align*}
    &\Gamma_{\lambda}(1) \\
    &= \sum_{\square_0~ \text{removable}} \Gamma_{\lambda \setminus \{\square_0\}}(1) \left( \frac{1-t^{\lambda_1-c(\square_0)+1}}{1-t} \right)\prod_{\square_0 \rightarrow \square} t^{d(\square)+1}\left(\frac{1-t^{c(\square)-c(\square_0)-1}}{1-t^{c(\square)-c(\square_0)+1}}\right) \cdots \left(\frac{1-t^{c(\square)-c(\square_0)-1+d(\square)}}{1-t^{c(\square)-c(\square_0)+1+d(\square)}}\right).\\
\end{align*}

Note that $\Gamma_{\emptyset}(1) = 1 \in k(t).$ Thus by induction on the size of $|\lambda|$ we see that $\Gamma_{\lambda}(1) \in k(t).$

\end{proof}

\end{document}
